\newcommand{\DoNotLoadEpstopdf}{}
\documentclass[11pt,letterpaper]{article}
\usepackage[T1]{fontenc}
\usepackage{lmodern,microtype}
\usepackage[margin=1in]{geometry}
\usepackage{amsmath,amssymb,amsthm,mathtools,booktabs,array,longtable}
\usepackage[dvipsnames]{xcolor}
\usepackage{enumitem,fancyhdr}
\usepackage[colorlinks=true,linkcolor=MidnightBlue,citecolor=MidnightBlue,urlcolor=MidnightBlue,breaklinks=true]{hyperref}
\hypersetup{pdftitle={Report 200: Beyond the leading count of commutative magmas},pdfauthor={Report 200},pdfsubject={Exact automorphism sectors, global remainders, rare symmetry and threshold inversion for OEIS A001425}}
\pdfinfoomitdate=1\pdftrailerid{}\pdfsuppressptexinfo=15
\pagestyle{fancy}\fancyhf{}
\fancyhead[L]{\small Commutative magmas: exponentially small symmetry}
\fancyhead[R]{\small Report 200}\fancyfoot[C]{\thepage}
\setlength{\headheight}{14pt}\setlength{\parskip}{4pt}\setlength{\parindent}{1.2em}
\emergencystretch=2em
\numberwithin{equation}{section}
\newtheorem{theorem}{Theorem}[section]
\newtheorem{proposition}[theorem]{Proposition}
\newtheorem{lemma}[theorem]{Lemma}
\newtheorem{corollary}[theorem]{Corollary}
\theoremstyle{definition}\newtheorem{definition}[theorem]{Definition}
\theoremstyle{remark}\newtheorem{remark}[theorem]{Remark}
\newcommand{\E}{\mathbb E}\newcommand{\Prob}{\mathbb P}
\newcommand{\Aut}{\operatorname{Aut}}\newcommand{\Fix}{\operatorname{Fix}}
\newcommand{\Sym}{\operatorname{Sym}}\newcommand{\ee}{\mathrm e}
\newcommand{\code}[1]{\texttt{\detokenize{#1}}}
% Added in the write (batch 112, 7 October 2026).
\newcommand{\file}[1]{\nolinkurl{#1}}
\newenvironment{writenote}[1][]{\par\smallskip\begingroup\small
 \noindent\textbf{[write]}\if\relax\detokenize{#1}\relax\else~\textbf{#1.}\fi\ }%
 {\par\endgroup\smallskip}
\title{\LARGE\bfseries Beyond the leading count\\of commutative magmas\\[9pt]\large Exact symmetry sectors, global remainders,\\rare automorphisms, and inverse thresholds}
\author{Report 200 \quad\textbullet\quad OEIS A001425}
\date{4 October 2026}
\begin{document}
\maketitle
\begin{abstract}
Let $U_n$ count isomorphism classes of arbitrary commutative binary operations on $n$ elements, with diagonal products included. The exact Burnside enumeration and the leading equivalent $U_n\sim n^{n(n+1)/2}/n!$ are already recorded in OEIS A001425. We study the much smaller symmetry excess after this identity contribution is removed. Every fixed moved cycle type has an exact contribution with a convergent, rational-coefficient amplitude in $1/n$. A uniform finite-$n$ bound controls all omitted conjugacy classes after any fixed defect, including large supports and fixed-point-free permutations. This gives an ordered hierarchy by defect, support, and polynomial power. We print the first transposition, $3$-cycle, and double-transposition sectors and derive labeled and unlabeled nonrigidity probabilities differing by a factor of two; conditional on nonrigidity, the automorphism group is asymptotically generated by one transposition. An exact-Gamma inverse retains the symmetry sectors and gives a one-sided-error two-ceiling enclosure. Its unrounded width is superexponentially small in the inverse-root variable, while the rounded enclosure can still contain two adjacent integers. Exact code, locally computed counts through $24$, and reproducible build instructions accompany the proofs. The conclusions are fixed-defect and eventual where stated; no effective inverse cutoff or absolute priority claim is made.
\end{abstract}
\begin{writenote}[Status and trust boundary, added 7 October 2026, batch~112]
This is bundle Report~200 of an external AI-assisted research session, filed
in ProveIt as the single-source research report
\file{a001425-commutative-magmas} (provenance, the sources as the write read
them, what was checked, relation to the repository, the collected non-claims
and the reading conventions are in Section~\ref{cmg:sec:provenance}). It is
unrefereed and not formalized. Its author line reads ``Report 200 \textbullet{}
OEIS A001425'' and its PDF author field ``Report 200'': it names no person or
tool. \emph{What rests on what.} Every theorem has a conventional proof in the
text: Burnside's lemma, an exact input-orbit count, an orbit-loss inequality
and a word-counting bound over all omitted permutations, analytic logarithms on
an explicit disk, and an orbit--stabilizer change of measure. No proof uses a
computation; the finite checks corroborate, and the decimals are diagnostics.
The write adds Remarks~\ref{cmg:rem:oeis} and~\ref{cmg:rem:transseries} and
dated notes; no statement, proof or number of the source was changed.
\end{writenote}
\clearpage
{\small\setcounter{tocdepth}{1}\tableofcontents}
\bigskip
\noindent\textbf{Reading guide.} Section~\ref{cmg:sec:main} states the scope and main result. Sections~\ref{cmg:sec:exact}--\ref{cmg:sec:amplitudes} prove the enumeration, global tail, and convergent amplitudes. Sections~\ref{cmg:sec:first}--\ref{cmg:sec:probability} extract the leading sectors and rare-symmetry laws. Section~\ref{cmg:sec:inverse} proves the inverse enclosure. The final sections describe exact verification, source boundaries, and questions that remain open.
\newpage

\section{The model and the scale of the question}\label{cmg:sec:main}
A \emph{commutative magma} on $[n]=\{1,\ldots,n\}$ is a binary operation satisfying $f(i,j)=f(j,i)$. We impose no associativity, identity element, cancellation, or idempotence. In particular $f(i,i)$ is a freely chosen table entry. In older enumeration terminology this is a commutative \emph{groupoid}; it is not a groupoid in the category-theoretic sense.

Equivalently, a labeled law is a map
\[
 f:\Sym^2([n])\longrightarrow[n],\qquad
 \Sym^2([n])=\{\{i,j\}:1\le i\le j\le n\}.
\]
Here pairs are unordered multisets, so the input set has
\begin{equation}\label{cmg:eq:N}
 N=N(n)=\frac{n(n+1)}2
\end{equation}
positions and there are exactly $n^N$ labeled laws. Let $U_n$ be the number of isomorphism classes. We set $U_0=1$, the count of the unique empty operation; normalized expressions involving $n^N$ will be used only for $n\ge1$. In exact products, $0^0=1$ and every empty product is $1$.

The symmetric group $S_n$ acts by relabeling input and output together. Write $\Fix_n(\sigma)$ for the number of laws for which $\sigma$ is an automorphism. Burnside's lemma yields
\begin{equation}\label{cmg:eq:burnside}
 U_n=\frac1{n!}\sum_{\sigma\in S_n}\Fix_n(\sigma),\qquad
 E_n:=\frac{n!U_n}{n^N}=\sum_{\sigma\in S_n}\frac{\Fix_n(\sigma)}{n^N}.
\end{equation}
Thus $E_n$ is also the mean automorphism-group order of a uniform labeled law. The identity contributes exactly $1$. The problem studied here is the excess $E_n-1$.

The exact count and
\begin{equation}\label{cmg:eq:known-leading}
 U_n\sim\frac{n^N}{n!}\sim\frac{\ee^n n^{n(n-1)/2}}{\sqrt{2\pi n}}
\end{equation}
are already in OEIS A001425 \cite{oeis}. An inverse-power Stirling expansion of $1/n!$ alone never sees $E_n-1$, which is smaller than every fixed inverse power of $n$. The additional structure below concerns the symmetry terms themselves.

\subsection{Fixed moved types}
Remove the fixed points of a permutation and record its remaining cycle lengths as a partition $\mu$ with all parts at least $2$. Write
\begin{equation}\label{cmg:eq:parameters}
 m_k=\#\{\text{parts of $\mu$ equal to $k$}\},\quad
 s=\sum_{k\ge2}km_k,\quad c=\sum_{k\ge2}m_k,\quad d=s-c,\quad
 z_\mu=\prod_{k\ge2}k^{m_k}m_k!.
\end{equation}
The integer $s$ is the moved support and $d$ the \emph{defect}, namely $n$ minus the total number of cycles including fixed points. For a nonempty type,
\begin{equation}\label{cmg:eq:support-range}
 d+1\le s\le2d.
\end{equation}
Let $T_\mu(n)$ denote its total contribution to the normalized sum \eqref{cmg:eq:burnside}. Its exact definition is given in \eqref{cmg:eq:sector}. It is convenient to set $T_\mu(n)=0$ when $s>n$, meaning that the type does not occur.

\begin{theorem}[Fixed-defect expansion with a global remainder]\label{cmg:thm:overview}
For a fixed integer $D\ge0$, put
\begin{equation}\label{cmg:eq:SD}
 S_D(n)=1+\sum_{1\le d(\mu)\le D}T_\mu(n),\qquad R_D(n)=E_n-S_D(n).
\end{equation}
Let $a=D+1$. For every integer $n\ge\max(8,4a)$,
\begin{equation}\label{cmg:eq:BD-main}
 0\le R_D(n)\le B_D(n),\qquad
 B_D(x)=\frac{x^{-ax+a(a+2)}}{1-x^{3-x/2}}+x^{x-3x^2/16}.
\end{equation}
In particular,
\begin{equation}\label{cmg:eq:Rorder}
 R_D(n)=O_D\!\left(n^{-(D+1)n+(D+1)(D+3)}\right).
\end{equation}
For each fixed nonempty type there are explicit rational numbers $\beta_\mu,C_\mu$ and a function $A_\mu$, analytic for $|t|<1/s$ with $A_\mu(0)=1$, such that for $n>s$,
\begin{equation}\label{cmg:eq:amplitude-main}
 T_\mu(n)=\frac{\ee^{C_\mu}}{z_\mu}\,
 n^{-dn+\beta_\mu}\ee^{-sn/2}A_\mu(1/n).
\end{equation}
All Taylor coefficients of $A_\mu$ are rational and are computable to any fixed order by an exact recurrence.
\end{theorem}

The bound is over the whole omitted permutation set. It is not obtained by comparing only the next few fixed-support types. The expansion orders sectors first by increasing $d$, then by increasing $s$, then by decreasing $\beta_\mu$; equal scales may be grouped. A finite algebraic truncation of an earlier amplitude generally leaves an error larger than a later exponential sector. To resolve that later sector, retain the earlier sector exactly or use its complete convergent amplitude. This distinction is used again in the inverse construction.

\subsection{Provenance, sources and reading conventions}\label{cmg:sec:provenance}
\begin{writenote}[Added 7 October 2026, batch~112]\raggedright
\emph{Provenance.} This report is Report~200 of the session bundle that
ProveIt's \file{docs/incoming} intake processed as batch~112: the archive
\file{Report200.zip} ($4{,}786{,}760$ bytes, 25 files, wrapper directory
\file{Report200/}), arrival commit \file{60f54ea06}, placed in
\file{f79c9bef1}. The text is the delivered \file{Report200.tex} (669 lines,
18 pages, dated 4~October 2026), shipped as \file{article.tex}. The programs
are shipped under \file{code/} (the root \file{build.py} beside them), their
usage note as \file{code-USAGE.txt}, the recorded outputs under \file{data/},
with \file{code/requirements.txt} as \file{data/code-requirements.txt} and the
source manifest (which also records the PDF-engine banner and the build epoch)
as \file{data/manifests-source_manifest.json}. Not shipped, and retrievable
from the arrival commit: the delivered \file{Report200.pdf}; the delivered
README, replaced by the report README; the two pure checksum manifests
\file{data/manifest.json} (11 entries) and
\file{manifests/package_manifest.json} (24 entries), verified at the write
together with the source manifest (10 entries); and the recorded output
\file{data/fixed_counts.json} ($4{,}150{,}875$ bytes, every cycle type for
$n\le24$ with its pair orbits and fixed counts), which \file{code/check.py}
regenerates (the report README gives the command; the write regenerated it from
the shipped files and found it equal to the delivered file after removing
carriage returns). The package records no ProveIt commit; it carries no
``prepared for private review'' line and no personal data. It is a
single-source report, so no merge choice was made. The write prefixed the 75
delivered labels with \file{cmg:} and updated the 41 references to them (31
\verb|\eqref|, 10 \verb|\ref|); it added this subsection, the notes marked
[write] and Remarks~\ref{cmg:rem:transseries} and~\ref{cmg:rem:oeis}, each the
last statement of its section, and inserted nothing before a delivered display
or statement, so every theorem, section and equation of the source keeps its
delivered number.

\emph{The sources, as the write read them} (7~October 2026): OEIS A001425 in
the internal format with its b-file (Remark~\ref{cmg:rem:oeis}); the
transseries volume
\path{Analysis/Transseries/docs/series-and-transseries/Transseries_And_Inversion/transseries_and_inversion.tex}
(\file{p0:def:model}, \file{p0:prop:factorial-core}, \file{p0:def:three-inverses},
\file{p0:thm:staircase}, \file{plt:def:lw-monomial-log-datum},
\file{plt:thm:lw-template}). Not read by the write: the cited literature
(Freese, Ture\v{c}ek, Burris, Harrison, Tamura); the source's account of it
(Section~\ref{cmg:sec:sources}) is reported, not confirmed.

\emph{What was checked at intake.} The batch-112 dossier read the manuscript,
found no demonstrably wrong claim, and reran \file{code/check.py} on a copy,
regenerating \file{data/fixed_counts.json}; the differences were Windows
carriage returns only (they also change the hashes in \file{data/manifest.json}).
\emph{What the write checked.} It read every proof and found no error. With
its own code (Python 3.14.4, exact integers and rationals; SymPy 1.14.0;
mpmath 1.3.0) it computed $U_n$ for $0\le n\le36$ by the grouped cycle-type
product, equal to all 37 terms of the OEIS b-file and to the 25 delivered
counts; checked the grouped product against a literal pair-orbit walk for
every cycle type with $n\le9$, and the five fixed counts printed in
Section~\ref{cmg:sec:exact}.1; enumerated all 729 laws with $n=3$ (the
automorphism distribution, the six statistics, $U_3=129$, $P_3=11/243$,
$Q_3=13/129$); recomputed the orbit data, $\beta_\mu$, $C_\mu$, $z_\mu$ of the
table in Section~\ref{cmg:sec:first} and, by a SymPy expansion of the exact
logarithms of~\eqref{cmg:eq:T2exact}--\eqref{cmg:eq:T22exact}, the
coefficients $L_1,\ldots,L_5$ of $T_2$ and the amplitude coefficients
of~\eqref{cmg:eq:T2exp}, \eqref{cmg:eq:T3exp} and~\eqref{cmg:eq:T22exp}; the
ratios $T_{22}/T_3=6.789890\ldots$, $1.304041\ldots$, $0.2078537\ldots$ at
$n=20$, $25$, $30$ (the printed values are roundings, as their notation
allows); $\sum_{n=1}^{24}p(n)=7337$; and $0\le R_D(n)\le B_D(n)$, with the true
bound~\eqref{cmg:eq:BD-main} at 50 digits, in all 62 cases $8\le n\le24$,
$0\le D<\lfloor n/4\rfloor$ (28 with odd $n$; the largest ratio
$R_D/B_D$ is $1.8\times10^{-4}$). Solving $F(\theta)=L$ with the exact Gamma
function at $L=10^3$, $10^5$, $10^8$, $10^{12}$, it found
$u\,(\theta-u-\frac12+5/(4\log u+2))$ between $0.64$ and $0.72$, as
Proposition~\ref{cmg:prop:initializer} predicts. It reran, on fresh copies,
\file{code/check.py} in normal and \texttt{-O} mode and from the shipped files
(every output equals the delivered one after removing carriage returns,
\file{data/manifest.json} apart) and \file{code/test_guards.py} (71 guard
checks, PASS). The builder was not run. Floating computations are
observations, not certificates.

\emph{Relation to the repository.} Before batch~112 no file of the repository
named A001425 or treated commutative magmas. The report
\file{a340021-distinguished-maximal-independent-sets} (batch~112) also
reduces an unlabeled count to its labeled count by Burnside's lemma, for a
different object; no result is shared, so no reciprocal note is proposed. No
Lean or Rocq file treats these sequences; no statement of this report is
formalized, and its place in the collection confers no formal status. The
inverse is compared with the transseries volume in
Remark~\ref{cmg:rem:transseries}.
\end{writenote}
\begin{writenote}[What is not claimed, collected from the source]
The exact Burnside count and the leading equivalent are in OEIS A001425; the
Burnside mechanism and the cycle-length product are background; Freese's
labeled/unlabeled framework and the event decomposition are classical context,
and his and Ture\v{c}ek's and Burris's unrestricted results are not applied to
the commutative space. The Harrison (1966) and Tamura (1970) originals and
Murski\u{\i}'s papers were not read, so no worldwide originality, absence of
an earlier proof or nonoverlap is claimed. The conclusions are fixed-defect and
eventual: no uniformity in growing defect or support, no optimal radius, no
effective inverse onset, no single-ceiling rule (the two ceilings can differ
arbitrarily far out), no subgroup-resolved higher rare-symmetry terms; a
finite Taylor truncation of an earlier sector does not resolve a later one;
floored-exponent bounds in the finite checks are test devices, not theorems;
numerical inverse and ordering values are diagnostics; hash manifests detect
alteration but do not authenticate. \emph{The write adds:} its checks are
exact or floating computations; its transseries remark claims no novelty for
any inversion.
\end{writenote}
\medskip
\begingroup\small
\noindent\emph{Reading conventions} (letters the source uses in more than one
sense, with the tempting false readings; no symbol was renamed).\par\smallskip
\renewcommand{\arraystretch}{1.1}
\noindent\begin{tabular}{@{}>{\raggedright\arraybackslash}p{0.8in}>{\raggedright\arraybackslash}p{5.4in}@{}}
\toprule
Symbol & Meanings\\
\midrule
$N$ & $N(n)=n(n+1)/2$ the number of input positions; $N_*(y)$ the integer threshold (the source distinguishes them).\\
$F$ & $F_\mu(n)$ the fixed-law count of~\eqref{cmg:eq:Fmu}; $F(x)$ the Gamma comparison function~\eqref{cmg:eq:F}. \emph{False reading:} $F(n)$ is $\log(n^N/n!)$, not a fixed-law count.\\
$E$, $\E$ & $E_n=n!U_n/n^N$, the mean automorphism-group order of a labeled law; $\E$ an expectation.\\
$P$, $Q$, $p$ & $P_n$ (labeled) and $Q_n\sim2P_n$ (unlabeled) nonrigidity probabilities; $p_j(\mu)$ the amplitude coefficients; $p_{\rm good}$. \emph{False reading:} the two probabilities differ by a factor two.\\
$H$, $h$ & $h_\ell(\mu)$ internal orbit counts; $H_\ell$ the orbit counts of a whole permutation in~\eqref{cmg:eq:mobius}; $H(x)=x^2\log x/2$; $H$ an invariant event.\\
$m$, $a$, $b$ & $m=n-s$ the fixed points and $m_k$ cycle multiplicities; $m_D(\theta)$ the derivative infimum. $a=D+1$; $a_\ell=s-b_\ell$; $a$ an automorphism order. $b_\ell$ the moved points fixed by $\sigma^\ell$.\\
$c$, $C$ & $c$ the number of moved cycles; $C_\mu$ the rational exponent constant.\\
$L$ & $L_r(\mu)$ the logarithmic amplitude coefficients; $L=\log y$ in Section~\ref{cmg:sec:inverse}.3.\\
$R$, $B$ & $R_D$ the global remainder, $R_1$, $R_2$; $B_D$ the explicit bound and $\underline B_D$ its floored version.\\
$U$, $u$ & $U_n$ the counts; $u$ the Lambert initializer.\\
$\mu$, $\mu_{\rm ar}$ & a moved cycle type; the arithmetic M\"obius function.\\
$\rho$ & a Cauchy radius $0<\rho<1/s$ in Section~\ref{cmg:sec:amplitudes}.\\
$X$, $Y$, $Z$ & numbers of transposition and other automorphisms; $Z=\{Y\ge1\}$ an event.\\
groupoid & the older word for magma used by OEIS, not a category-theoretic groupoid (the source says so).\\
\bottomrule
\end{tabular}
\endgroup

\section{Exact enumeration by equivariant maps}\label{cmg:sec:exact}
\subsection{The orbit formula}
The automorphism condition is
\begin{equation}\label{cmg:eq:equivariant}
 f(\{\sigma i,\sigma j\})=\sigma f(\{i,j\}).
\end{equation}
Suppose an input orbit under $\sigma$ has length $\ell$. Choosing the output $y$ at one representative forces every other output along that orbit. The assignment closes consistently if and only if $\sigma^\ell y=y$. Hence it has exactly $\Fix(\sigma^\ell)$ choices, independently of all other input orbits.

Let $h_\ell(\mu)$ be the number of length-$\ell$ orbits on $\Sym^2([s])$ for a permutation with moved type $\mu$. These are the \emph{internal} input orbits, with diagonal pairs included. Set
\begin{equation}\label{cmg:eq:bq}
 b_\ell=\sum_{k\mid\ell}km_k,\qquad q=\sum_{\ell\ge1}h_\ell,\qquad m=n-s.
\end{equation}
Thus $b_\ell$ is the number of moved points fixed by the $\ell$th power, and $\Fix(\sigma^\ell)=m+b_\ell$.

\begin{proposition}[Exact sector]\label{cmg:prop:exact}
For $n\ge s$, the fixed-law count and sector are
\begin{align}
 F_\mu(n)&=m^{m(m+1)/2}
 \prod_{k\ge2}(m+b_k)^{m m_k}
 \prod_{\ell\ge1}(m+b_\ell)^{h_\ell},\label{cmg:eq:Fmu}\\
 T_\mu(n)&=\frac{(n)_s}{z_\mu n^{N(n)}}F_\mu(n),
 \qquad (n)_s=n(n-1)\cdots(n-s+1),\label{cmg:eq:sector}\\
 E_n&=1+\sum_{\substack{\mu\ne\varnothing\\s(\mu)\le n}}T_\mu(n).\label{cmg:eq:Esector}
\end{align}
Only factors with nonzero $m_k$ or $h_\ell$ need to be included.
\end{proposition}
\begin{proof}
The fixed--fixed pairs give $m(m+1)/2$ singleton input orbits and $m$ choices per orbit. Crossing one fixed point with a moved $k$-cycle gives one input orbit of length $k$, so there are $m m_k$ such orbits. The remaining input orbits are precisely those counted by $h_\ell$. The preceding equivariance argument proves \eqref{cmg:eq:Fmu}. There are $n!/((n-s)!z_\mu)=(n)_s/z_\mu$ permutations with this moved type. Grouping \eqref{cmg:eq:burnside} by type proves the result.
\end{proof}

The convention $0^0=1$ matters when $m=0$. The fixed--fixed exponent and all mixed exponents then vanish, but an internal orbit with no admissible output gives zero. There is no division by $m$. For example a fixed-point-free transposition has one singleton input orbit $\{1,2\}$, so $F_{(2)}(2)=0$. A fixed-point-free $3$-cycle instead has two length-$3$ input orbits, so $F_{(3)}(3)=3^2=9$. Further examples are $F_{(4)}(4)=0$, $F_{(3,3)}(6)=6^7$, and $F_{(3,6)}(9)=3^3\cdot9^6$. Fixed-point-free cases are therefore neither all zero nor all governed by one parity rule.

\subsection{Two ways to compute the input orbits}
A literal orbit walk on the $s(s+1)/2$ unordered pairs computes all $h_\ell$. There is also a trace-and-inversion route that does not walk these pair orbits. For any permutation $\tau$, let $f_r=\Fix(\tau^r)$ and $H_\ell$ be the number of length-$\ell$ orbits on its entire unordered-pair input set. Then
\begin{equation}\label{cmg:eq:trace}
 \Fix(\Sym^2(\tau^r))=\frac{f_r^2+f_{2r}}2
   =\sum_{\ell\mid r}\ell H_\ell.
\end{equation}
Indeed fixed input pairs are either drawn from individually fixed points, or are the endpoints of a $2$-cycle of $\tau^r$. There are $f_r(f_r+1)/2$ of the former and $(f_{2r}-f_r)/2$ of the latter. If $\mu_{\rm ar}$ is the arithmetic M\"obius function, inversion gives
\begin{equation}\label{cmg:eq:mobius}
 H_\ell=\frac1\ell\sum_{r\mid\ell}\mu_{\rm ar}(\ell/r)\,
 \frac{f_r^2+f_{2r}}2,\qquad
 \Fix_n(\tau)=\prod_{\ell\ge1}f_\ell^{H_\ell}.
\end{equation}
Only divisors of the order of $\tau$ can occur as input-orbit lengths. Formula \eqref{cmg:eq:mobius} gives an independent exact computational check of \eqref{cmg:eq:Fmu}.

For comparison with the cycle-length formula in OEIS, inside a single $k$-cycle there are $(k+1)/2$ length-$k$ orbits if $k$ is odd. If $k$ is even, there are $k/2$ length-$k$ orbits and one length-$k/2$ antipodal orbit. Between distinct $k$- and $j$-cycles there are $\gcd(k,j)$ orbits of length $\operatorname{lcm}(k,j)$. Multiplying the appropriate fixed-output counts is precisely the classical grouped product used by the entry, including its half-length factor for even cycles. The Burnside mechanism is background, not a claim of novelty.

\section{A tail bound over every omitted permutation}\label{cmg:sec:tail}
Let $O(\sigma)$ be the total number of input orbits on $\Sym^2([n])$. Each such orbit has at most $n$ admissible outputs, so
\begin{equation}\label{cmg:eq:orbit-bound}
 \Fix_n(\sigma)\le n^{O(\sigma)}.
\end{equation}
This holds also when the fixed count is zero.

\begin{lemma}[Orbit loss]\label{cmg:lem:loss}
For a permutation of moved support $s$ and positive defect $d$,
\begin{align}
 N-O(\sigma)&\ge(n-s)d+s^2/4,\label{cmg:eq:loss-support}\\
 N-O(\sigma)&\ge
 \begin{cases}d(n-d),&d\le n/2,\\ n^2/4,&d\ge n/2.\end{cases}\label{cmg:eq:loss-defect}
\end{align}
\end{lemma}
\begin{proof}
The fixed--moved pairs number $(n-s)s$ and form $(n-s)c$ orbits. Their loss is exactly $(n-s)(s-c)=(n-s)d$. On the moved support there are $M=s(s+1)/2$ pairs. Exactly $m_2$ are fixed inputs: the endpoints of the individual transpositions. A diagonal cannot be fixed because its point is moved. Every other internal orbit has length at least two. Consequently
\[
 q\le m_2+\frac{M-m_2}{2}=\frac{M+m_2}{2}
 \le\frac{s(s+2)}4,
\]
and $M-q\ge s^2/4$. This proves \eqref{cmg:eq:loss-support}.

For fixed $d$, the right side of \eqref{cmg:eq:loss-support} has derivative $-d+s/2$ in $s$. It is nonincreasing on the enlarged interval $0\le s\le\min(2d,n)$, which contains every possible support. At its right endpoint it equals $d(n-d)$ if $2d\le n$ and $n^2/4$ otherwise. Minimizing on this larger interval is in the safe direction and proves \eqref{cmg:eq:loss-defect}.
\end{proof}

\begin{lemma}[Counting by defect]\label{cmg:lem:defect-count}
The number of permutations in $S_n$ with defect $d$ is at most $n^{2d}$.
\end{lemma}
\begin{proof}
A cycle of length $k$ is a product of $k-1$ transpositions. Decomposing all cycles shows that a defect-$d$ permutation is the value of some word of exactly $d$ transpositions. There are $\binom n2^d\le n^{2d}$ such words. This is a surjective encoding bound; uniqueness is neither asserted nor needed.
\end{proof}

\begin{proof}[Proof of the remainder part of Theorem~\ref{cmg:thm:overview}]
Put $a=D+1$ and suppose $n\ge\max(8,4a)$. For $a\le d\le\lfloor n/4\rfloor$, Lemmas~\ref{cmg:lem:loss} and~\ref{cmg:lem:defect-count} bound the total normalized contribution at defect $d$ by
\[
 H_d=n^{2d-d(n-d)}.
\]
Consecutive majorants have ratio
\[
 \frac{H_{d+1}}{H_d}=n^{2d+3-n}\le n^{3-n/2}<1.
\]
The initial range is nonempty, even at the boundary $n=4a$. Extending its finite sum to a geometric series gives
\begin{equation}\label{cmg:eq:small-defect-tail}
 \sum_{a\le d\le\lfloor n/4\rfloor}\ \sum_{\mu:d(\mu)=d}T_\mu(n)
 \le\frac{n^{-an+a(a+2)}}{1-n^{3-n/2}}.
\end{equation}

For all remaining defects $d>n/4$, the loss in \eqref{cmg:eq:loss-defect} is at least $3n^2/16$: on $[n/4,n/2]$ the function $d(n-d)$ is increasing, and beyond $n/2$ the bound $n^2/4$ is stronger. There are at most $n!\le n^n$ permutations in this entire second range. Thus its total normalized contribution is at most $n^{n-3n^2/16}$. These two ranges cover every omitted type, including macroscopic support and fixed-point-free permutations. Their sum is $B_D(n)$.

For fixed $a$, the denominator in the first summand tends to $1$. The second summand is negligible compared with the first, since its logarithm has negative quadratic rather than negative linear leading order in $n$, multiplied by $\log n$. This proves \eqref{cmg:eq:Rorder}.
\end{proof}

\begin{proposition}[Monotonicity of the explicit bound]\label{cmg:prop:Bdecreasing}
The function $B_D(x)$ is strictly decreasing on $x\ge\max(8,4(D+1))$.
\end{proposition}
\begin{proof}
For $a=D+1$, put $A(x)=x^{-ax+a(a+2)}$ and $r(x)=x^{3-x/2}$. Their logarithmic derivatives are
\[
 (\log A)'=a\left(-\log x-1+\frac{a+2}{x}\right)<0,
 \qquad
 (\log r)'=-\frac{\log x}{2}-\frac12+\frac3x<0.
\]
Here $(a+2)/x\le1/2$ if $a\ge2$, and is at most $3/8$ if $a=1$. Since $0<r<1$, $A/(1-r)$ is strictly decreasing. The logarithmic derivative of the other summand is
\[
 (1-3x/8)\log x+1-3x/16<0\qquad(x\ge8).
\]
Both positive summands decrease.
\end{proof}

\section{Convergent amplitudes and exact coefficients}\label{cmg:sec:amplitudes}
For the internal orbit counts of \eqref{cmg:eq:bq}, define
\begin{equation}\label{cmg:eq:betaC}
 \beta_\mu=s-cs+q+\frac{s(s-1)}2,\qquad
 C_\mu=\frac{3s^2}4-\frac s2-cs+\sum_{k\ge2}m_k b_k.
\end{equation}
The integer $\beta_\mu$ is the polynomial exponent of the sector. The constant $C_\mu$ is rational; the exponential leading constant is $\ee^{C_\mu}/z_\mu$.

\begin{theorem}[Analytic fixed-sector factorization]\label{cmg:thm:analytic}
For every fixed nonempty moved type $\mu$, \eqref{cmg:eq:amplitude-main} holds for $n>s$, with $A_\mu$ holomorphic in $|t|<1/s$ and $A_\mu(0)=1$. Write
\[
 \log A_\mu(t)=\sum_{r\ge1}L_r(\mu)t^r,\qquad
 A_\mu(t)=\sum_{j\ge0}p_j(\mu)t^j.
\]
Then
\begin{align}\label{cmg:eq:Lr}
 L_r={}&-\frac1r\sum_{j=0}^{s-1}j^r
 -\frac{s^{r+2}}{2(r+2)}
 -\frac{(1-2s)s^{r+1}}{2(r+1)}
 -\frac{(s^2-s)s^r}{2r}\notag\\
 &+\sum_{k\ge2}m_k\left[-\frac{(s-b_k)^{r+1}}{r+1}
                         +\frac{s(s-b_k)^r}{r}\right]
 -\sum_{\ell\ge1}h_\ell\frac{(s-b_\ell)^r}{r},\\
 p_0={}&1,\qquad
 p_j=\frac1j\sum_{r=1}^j rL_r p_{j-r}\quad(j\ge1).\label{cmg:eq:p-rec}
\end{align}
In particular all $L_r$ and $p_j$ are rational.
\end{theorem}
\begin{proof}
The total number of input orbits in \eqref{cmg:eq:Fmu} is
\[
 K=\frac{(n-s)(n-s+1)}2+(n-s)c+q.
\]
Factoring a power of $n$ from every base and $n^s$ from $(n)_s$ gives
\[
 s+K-N=-dn+\beta_\mu.
\]
Put $a_\ell=s-b_\ell$ and $t=1/n$. The remaining logarithm before removal of the pole and constant is
\begin{align}\label{cmg:eq:rawlog}
 J_\mu(t)={}&\sum_{j=0}^{s-1}\log(1-jt)
 +\frac{(t^{-1}-s)(t^{-1}-s+1)}2\log(1-st)\notag\\
 &+(t^{-1}-s)\sum_{k\ge2}m_k\log(1-a_k t)
 +\sum_{\ell\ge1}h_\ell\log(1-a_\ell t).
\end{align}
Every logarithm is taken on the branch analytic at zero. Since $0\le a_\ell\le s$ and $0\le j<s$, these logarithms are analytic on $|t|<1/s$. Their Laurent combination has only one negative-power term, $-s/(2t)$. Expanding through degree zero shows that its constant is
\[
 -\frac{s^2}4-\frac{(1-2s)s}{2}-\sum_km_k(s-b_k)=C_\mu.
\]
Therefore
\begin{equation}\label{cmg:eq:Adefinition}
 A_\mu(t)=\exp\left(J_\mu(t)+\frac{s}{2t}-C_\mu\right)
\end{equation}
is analytic after the removable singularity at zero is filled in, and its value there is $1$. This proves the exact factorization and convergence. Collecting the coefficient of $t^r$ in \eqref{cmg:eq:rawlog} gives \eqref{cmg:eq:Lr}. Finally differentiate $A_\mu'=A_\mu(\log A_\mu)'$ and equate coefficients to obtain \eqref{cmg:eq:p-rec}.
\end{proof}

For a concrete fixed-order bound, choose $0<\rho<1/s$ and put $M_\rho=\max_{|t|=\rho}|A_\mu(t)|$. Cauchy's coefficient bound gives, for $n>1/\rho$,
\begin{equation}\label{cmg:eq:Cauchy}
 \left|A_\mu(1/n)-\sum_{j=0}^J\frac{p_j}{n^j}\right|
 \le \frac{M_\rho}{(n\rho)^{J+1}(1-1/(n\rho))}.
\end{equation}
The result is not merely a formal series. This bound is for fixed $\mu$; neither uniformity for growing support nor an optimal radius is asserted.

\subsection{Ordering the sectors}
The factorization implies
\begin{equation}\label{cmg:eq:logscale}
 \log T_\mu(n)=-dn\log n-\frac{s n}{2}
              +\beta_\mu\log n+C_\mu-\log z_\mu+O_\mu(n^{-1}).
\end{equation}
For two fixed types, a smaller defect wins over every fixed difference in support because $n\log n$ dominates $n$. At equal defect, the smaller support wins over every fixed polynomial factor. At equal defect and support, the larger $\beta$ wins polynomially. If all three agree, the leading constants can be added, and further coefficients compared if needed.

There are finitely many types at each defect: subtracting $1$ from every cycle length turns the type into an integer partition of $d$. The number through defect $4$ is $1+2+3+5=11$. For every retained type with $d\le D$ and every fixed $M\ge0$,
\begin{equation}\label{cmg:eq:tail-flat}
 \frac{n^M R_D(n)}{T_\mu(n)}\longrightarrow0.
\end{equation}
Indeed the logarithm of this ratio is at most $-(D+1-d)n\log n+sn/2+O_{D,\mu,M}(\log n)$, which tends to $-\infty$. This proves separation from the full global tail.

\begin{remark}[What an all-orders sector expansion means]\label{cmg:rem:all-orders}
For any fixed $D$, the finite family of exact sectors resolves the remainder scale \eqref{cmg:eq:Rorder}; within each sector, its amplitude has a convergent expansion with coefficients to arbitrary order. It does not follow that one can replace every retained amplitude by the same finite Taylor polynomial while keeping an exponentially small final remainder. A nonzero next algebraic coefficient in an earlier sector produces an error on that earlier exponential scale. The safe representation is the exact sector sum plus $R_D$, with Taylor expansions used to describe the individual sectors.
\end{remark}

\section{The first sectors explicitly}\label{cmg:sec:first}
For readability write $T_2=T_{(2)}$, $T_3=T_{(3)}$, and $T_{22}=T_{(2,2)}$. The first internal orbit data are
\[
\begin{array}{c|c|r|r|r|r|r}
 \mu&(h_\ell)_{h_\ell>0}&d&s&\beta&C&z\\\hline
 (2)&h_1=1,\ h_2=1&1&2&3&2&2\\
 (3)&h_3=2&2&3&5&21/4&3\\
 (2,2)&h_1=2,\ h_2=4&2&4&8&10&8
\end{array}
\]
For the transposition, the one singleton internal pair is its two endpoints; its two diagonals form one length-$2$ orbit. For two transpositions there are two singleton endpoint pairs and four length-$2$ internal orbits.

\subsection{A single transposition}
The exact expression, valid for $n\ge2$ under the zero-power convention, is
\begin{equation}\label{cmg:eq:T2exact}
 T_2(n)=\frac{n(n-1)}2\,
 \frac{(n-2)^{(n-2)(n-1)/2+1}\,n^{n-1}}{n^{n(n+1)/2}}.
\end{equation}
For $n>2$ it may also be written
\[
 T_2(n)=\frac12\left(1-\frac1n\right)n^{3-n}
        \left(1-\frac2n\right)^{(n^2-3n+4)/2}.
\]
The first logarithmic amplitude coefficients are $-10/3,-5/2,-43/15,-239/60,-43/7$. Exponentiating gives
\begin{align}\label{cmg:eq:T2exp}
 T_2(n)=\frac{\ee^2}2\ee^{-n}n^{3-n}\biggl[&1-\frac{10}{3n}
 +\frac{55}{18n^2}-\frac{286}{405n^3}
 -\frac{463}{9720n^4}\notag\\[-2pt]
 &-\frac{781}{20412n^5}+O(n^{-6})\biggr].
\end{align}
This is the leading symmetry excess, not a correction coming from $n!$.

\subsection{A $3$-cycle and two transpositions}
The next two sectors in eventual asymptotic order are
\begin{align}
 T_3(n)&=\frac{(n)_3}{3}\,
 \frac{(n-3)^{(n-3)(n-2)/2}\,n^{n-1}}{n^{n(n+1)/2}},\label{cmg:eq:T3exact}\\
 T_{22}(n)&=\frac{(n)_4}{8}\,
 \frac{(n-4)^{(n-4)(n-3)/2+2}\,n^{2n-4}}{n^{n(n+1)/2}}.\label{cmg:eq:T22exact}
\end{align}
Their amplitude expansions begin
\begin{align}
 T_3(n)&=\frac{\ee^{21/4}}3\ee^{-3n/2}n^{5-2n}
 \left[1-\frac{21}{4n}+\frac{325}{32n^2}-\frac{5607}{640n^3}+O(n^{-4})\right],\label{cmg:eq:T3exp}\\
 T_{22}(n)&=\frac{\ee^{10}}8\ee^{-2n}n^{8-2n}
 \left[1-\frac{62}{3n}+\frac{1667}{9n^2}-\frac{383402}{405n^3}+O(n^{-4})\right].\label{cmg:eq:T22exp}
\end{align}
Thus, for every $n\ge12$,
\begin{equation}\label{cmg:eq:three-sector}
 E_n=1+T_2(n)+T_3(n)+T_{22}(n)+R_2(n),\quad
 0\le R_2(n)\le\frac{n^{-3n+15}}{1-n^{3-n/2}}+n^{n-3n^2/16}.
\end{equation}
The $T$ terms in this equality are the exact expressions. The displayed finite algebraic expansions do not replace them in the certified remainder claim.

Ordering is eventual rather than a small-$n$ numerical rule. In particular,
\[
 \frac{T_{22}}{T_3}\sim\frac38\ee^{19/4}\,n^3\ee^{-n/2}\longrightarrow0.
\]
Yet this ratio is about $6.78989$ at $n=20$, $1.30404$ at $n=25$, and $0.207854$ at $n=30$. The polynomial factor can delay the eventual support ordering.

\subsection{Beyond defect two}
At defects $3$ and $4$ the moved types are respectively
\begin{align*}
 d=3:&\quad (4),\ (2,3),\ (2,2,2),\\
 d=4:&\quad (5),\ (2,4),\ (3,3),\ (2,2,3),\ (2,2,2,2).
\end{align*}
The exact coefficient recurrence treats all of them without a separate asymptotic argument. The companion data give the parameters, internal orbit counts, and five correction coefficients for all eleven types through defect $4$. At equal support, different types need not have the same polynomial exponent; therefore support alone is not the last level in the hierarchy.

\section{Rare symmetry in the two probability models}\label{cmg:sec:probability}
A law is \emph{rigid} if its automorphism group is trivial. Let $P_n$ be the nonrigidity probability for a uniformly chosen labeled law. Let $Q_n$ be the nonrigidity probability for a uniformly chosen isomorphism class. These are different probability spaces, although both probabilities tend to zero.

\begin{theorem}[Nonrigidity and its typical group]\label{cmg:thm:prob}
As $n\to\infty$,
\begin{equation}\label{cmg:eq:PQ}
 P_n\sim T_2(n),\qquad Q_n\sim2T_2(n).
\end{equation}
In either model, conditional on nonrigidity, the probability that the automorphism group is generated by a single transposition tends to $1$.
\end{theorem}
\begin{proof}
For a uniform labeled law, let $X$ count transposition automorphisms, and let $Y$ count all other nonidentity automorphisms. Then exactly
\begin{equation}\label{cmg:eq:AXY}
 |\Aut|=1+X+Y,\qquad \E X=T_2,\qquad \E Y=R_1.
\end{equation}
The last equality holds because the transpositions are precisely the permutations of defect $1$. Define $Z=\{Y\ge1\}$, the event that a nonidentity nontransposition occurs. The first-moment bound gives $\Prob(Z)\le R_1$.

The product of two distinct transpositions is a $3$-cycle if they overlap and a double transposition if they are disjoint. In either case it is a nonidentity nontransposition. Hence outside $Z$, at most one transposition can occur. Write $p_{\rm good}$ for the probability that $\Aut=\{1,\tau\}$ for a single transposition $\tau$. Then
\begin{align}
 0\le T_2-p_{\rm good}&=\E[X;Z]\le\binom n2R_1,\label{cmg:eq:good}\\
 P_n&=p_{\rm good}+\Prob(Z),\qquad
 |P_n-T_2|\le\left(\binom n2+1\right)R_1.\label{cmg:eq:Pbound}
\end{align}
By the tail theorem, $R_1=O(n^{-2n+8})$. Combining this with \eqref{cmg:eq:T2exp},
\[
 \frac{n^2R_1}{T_2}=O\!\left(\ee^n n^{-n+7}\right)\longrightarrow0.
\]
This proves the labeled assertions, including the conditional group statement.

For any isomorphism-invariant event $H$, orbit--stabilizer gives the exact change of measure
\begin{equation}\label{cmg:eq:measure}
 \Prob_{\rm unlabeled}(H)=\frac{\E_{\rm labeled}[|\Aut|\,\mathbf1_H]}{E_n}.
\end{equation}
Indeed a class with automorphism order $a$ has $n!/a$ labelings, which are weighted by $a$ in the numerator. For rigidity, where $|\Aut|=1$, this gives
\begin{equation}\label{cmg:eq:Qexact}
 Q_n=1-\frac{1-P_n}{E_n}=\frac{P_n+E_n-1}{E_n}.
\end{equation}
Since $E_n-1=T_2+R_1\sim T_2$ and $E_n\to1$, the unlabeled equivalent is $2T_2$.

Finally, $Y$ vanishes outside $Z$, so
\begin{equation}\label{cmg:eq:exceptional}
 \E[|\Aut|;Z]=\Prob(Z)+\E[X;Z]+\E Y
 \le\left(\binom n2+2\right)R_1.
\end{equation}
After division by $E_n$ this is negligible relative to $Q_n\sim2T_2$. The conditional unlabeled group statement follows.
\end{proof}

The proof uses the event of containing a nontransposition, rather than an inference from group order alone. A subgroup of order two can be generated by a product of several disjoint transpositions, so order two by itself does not identify the permutation type. The event decomposition is consistent with the classical framework of Freese \cite{freese}, but the estimates here are proved for the commutative table space itself.

Higher-order symmetry probabilities cannot be read off by dividing individual Burnside sectors by a putative typical group order. Different automorphisms can belong to the same law; their fixed-law sets intersect according to the subgroups they generate. The first-order result avoids this issue by showing that the entire exceptional mass is negligible. Subgroup-resolved higher corrections remain a separate problem.

\section{Inverse thresholds with exact identity background}\label{cmg:sec:inverse}
\subsection{The smooth comparison functions}
To invert the rapid growth without losing the exponentially small symmetry scale, retain the identity term exactly through the Gamma function:
\begin{equation}\label{cmg:eq:F}
 F(x)=\frac{x(x+1)}2\log x-\log\Gamma(x+1),\qquad x>0.
\end{equation}
At integers, $F(n)=\log(n^{N(n)}/n!)$. For a fixed $D\ge0$ and real $x>2D$, extend each retained sector using \eqref{cmg:eq:Fmu}--\eqref{cmg:eq:sector} with positive real bases and real powers, and put
\begin{equation}\label{cmg:eq:GD}
 G_D(x)=F(x)+\log\left(1+\sum_{1\le d(\mu)\le D}T_\mu(x)\right).
\end{equation}
For $D=0$ this means $G_0=F$. All factors are positive on the stated domain because $s\le2D<x$.

\begin{proposition}[Eventual monotonicity and one-sided error]\label{cmg:prop:inverse-error}
For each fixed $D$, $G_D'(x)\sim x\log x$. The function $G_D$ is eventually strictly increasing and unbounded. For every integer $n\ge\max(8,4(D+1))$,
\begin{equation}\label{cmg:eq:logerror}
 0\le\log U_n-G_D(n)\le B_D(n).
\end{equation}
The integer sequence $U_n$ is eventually strictly increasing.
\end{proposition}
\begin{proof}
Differentiating a fixed-sector factorization on the positive real axis gives
\[
 \frac{T_\mu'(x)}{T_\mu(x)}=-d\log x-d-\frac s2+O_\mu(x^{-1}).
\]
The convergent amplitude justifies the differentiated error: its contribution is $O_\mu(x^{-2})$. The finite sector sum and its derivative are therefore negligible relative to the leading identity derivative. The standard Gamma asymptotic gives
\[
 F'(x)=x\log x+\frac x2-\frac{\log x}{2}+\frac12+O(x^{-1}),
\]
so $G_D'(x)\sim x\log x$.

At integers, \eqref{cmg:eq:SD} and \eqref{cmg:eq:burnside} give exactly
\[
 \log U_n-G_D(n)=\log\left(1+\frac{R_D(n)}{S_D(n)}\right).
\]
As $S_D(n)\ge1$, the inequalities $0\le\log(1+u)\le u$ prove \eqref{cmg:eq:logerror}. Finally $G_D(n+1)-G_D(n)\to\infty$ while $B_D(n)\to0$, so eventually
\[
 \log U_{n+1}\ge G_D(n+1)>G_D(n)+B_D(n)\ge\log U_n.
\]
This proves eventual strict monotonicity rather than assuming it from the computed terms.
\end{proof}

Replacing $F$ by a finite Stirling truncation introduces a logarithmic error of algebraic size, much larger than the symmetry corrections. Likewise, replacing an earlier sector in $G_D$ by a finite amplitude polynomial generally destroys the precision needed to resolve a later sector. Stirling's formula is useful for initial guesses; the certified comparison function \eqref{cmg:eq:GD} keeps both the exact Gamma term and the exact finite-defect sectors.

\subsection{A two-ceiling enclosure}
Define the integer threshold
\begin{equation}\label{cmg:eq:Nthreshold}
 N_*(y)=\min\{n\ge0:U_n\ge y\}.
\end{equation}
The notation distinguishes this inverse from the input count $N(n)$. For sufficiently large $y$, let $\theta=\theta_D(y)$ be the unique solution on the eventual increasing branch of
\begin{equation}\label{cmg:eq:theta}
 G_D(\theta)=\log y.
\end{equation}
Set
\begin{equation}\label{cmg:eq:eta}
 m_D(\theta)=\inf_{t\in[\theta-1,\theta]}G_D'(t),\qquad
 \eta_D(\theta)=\frac{B_D(\theta-1)}{m_D(\theta)}.
\end{equation}
For large $\theta$, the denominator is positive and is asymptotic to $\theta\log\theta$.

\begin{theorem}[One-sided-error inverse enclosure]\label{cmg:thm:inverse}
For every fixed $D\ge0$ and all sufficiently large targets $y$,
\begin{equation}\label{cmg:eq:ceilings}
 \left\lceil\theta_D(y)-\eta_D(\theta_D(y))\right\rceil
 \le N_*(y)\le\left\lceil\theta_D(y)\right\rceil.
\end{equation}
For $a=D+1$,
\begin{equation}\label{cmg:eq:eta-asymptotic}
 \eta_D(\theta)\sim
 \frac{\ee^a\theta^{-a\theta+a(a+3)-1}}{\log\theta},\qquad
 \log\eta_D(\theta)=-a\theta\log\theta+O_D(\log\theta).
\end{equation}
The underlying real interval has width $\eta_D(\theta)$, which decays faster than every fixed exponential in $\theta$. After taking ceilings the enclosure has width at most one for all sufficiently large targets; it need not always be a singleton.
\end{theorem}
\begin{proof}
A finite exceptional prefix of $U_n$ is irrelevant once $y$ exceeds its maximum. All integer comparisons below may therefore be made in the eventual range. Since $G_D$ is increasing there and $\log U_n\ge G_D(n)$, the integer $\lceil\theta\rceil$ reaches the target, proving the upper bound.

Eventually $G_D(n)+B_D(n)<G_D(n+1)$. Thus any candidate integer $n<\theta-1$ satisfies
\[
 \log U_n\le G_D(n)+B_D(n)<G_D(n+1)<G_D(\theta)=\log y
\]
and cannot reach the target. If instead $\theta-1\le n<\theta$ and $U_n\ge y$, integration of the derivative lower bound gives
\[
 m_D(\theta)(\theta-n)\le G_D(\theta)-G_D(n)
 \le B_D(n)\le B_D(\theta-1),
\]
where the last inequality follows from Proposition~\ref{cmg:prop:Bdecreasing}. Hence $n\ge\theta-\eta_D(\theta)$, proving the lower ceiling.

The first summand of $B_D(\theta-1)$ is asymptotic to
\[
 (\theta-1)^{-a\theta+a(a+3)}
 \sim\ee^a\theta^{-a\theta+a(a+3)}.
\]
Its denominator tends to $1$ and the second summand is negligible. Dividing by $m_D(\theta)\sim\theta\log\theta$ proves \eqref{cmg:eq:eta-asymptotic}. In particular $0<\eta_D(\theta)<1$ eventually, so the ceilings differ by at most one.
\end{proof}

The shrinking scale in this theorem is $\theta$, not the target $y$. Nor does a shrinking real width imply eventual exact rounding: for every sufficiently large integer $k$, choose $\theta>k$ close enough that $\theta-k<\eta_D(\theta)<1$. The two ceilings are then $k$ and $k+1$. In fact the upper ceiling alone can give the wrong threshold arbitrarily far out. An omitted single-cycle type of length $D+2$ has positive contribution for all sufficiently large $k$, so
\[
 U_k>\exp G_D(k).
\]
Targets just above $\exp G_D(k)$ and at most $U_k$ have true threshold $k$ but smooth-root upper ceiling $k+1$, once the earlier terms are below the target. The eventual gap estimate ensures that last condition. No single-ceiling rule is asserted.

No explicit numerical onset for the eventual statements is proved here. A numerical evaluation of $\theta$ and $\eta$ by ordinary floating arithmetic is also not automatically an interval certificate; a certified implementation must control the root, derivative infimum, and bound with directed error estimates.

\subsection{A Lambert-$W$ initializer}
Let $L=\log y$ and let $W$ be the positive real Lambert function. Put
\begin{equation}\label{cmg:eq:u}
 u=\sqrt{\frac{4L}{W(4L)}}.
\end{equation}
Then $H(u)=L$ for $H(x)=x^2\log x/2$: writing $w=W(4L)$ gives $u=\ee^{w/2}$ and $u^2\log u/2=w\ee^w/4=L$.

\begin{proposition}[First explicit inverse center]\label{cmg:prop:initializer}
For every fixed $D$,
\begin{equation}\label{cmg:eq:initializer}
 \theta_D(y)=u+\frac{\frac12\log u-1}{\log u+\frac12}+O_D(u^{-1})
 =u+\frac12-\frac{5}{4\log u+2}+O_D(u^{-1}).
\end{equation}
\end{proposition}
\begin{proof}
Stirling's expansion, used only for this initializer, gives
\begin{equation}\label{cmg:eq:Fstirling}
 F(x)=H(x)-\frac{x\log x}{2}+x-\frac{\log x}{2}
 -\frac{\log(2\pi)}2-\frac1{12x}+O(x^{-3}).
\end{equation}
Let $\delta=(\frac12\log u-1)/(\log u+\frac12)$. Since $H'(u)=u(\log u+\frac12)$, the contribution $H'(u)\delta$ cancels the $-u\log u/2+u$ discrepancy. With $\delta=O(1)$, Taylor's formula leaves $F(u+\delta)-L=O(\log u)$. The finite-sector logarithm in $G_D-F$ is smaller than every power of $u^{-1}$, so the same residual estimate holds for $G_D$. Division by $G_D'\sim u\log u$, with the mean-value theorem on a bounded neighborhood, gives root error $O_D(u^{-1})$. Elementary algebra gives the second displayed expression.
\end{proof}

Higher algebraic inverse terms can be generated by further Stirling and Newton expansions. They accelerate initialization; they cannot replace the exact comparison function when exponentially small symmetry information is required.

\begin{remark}[{[write]} The inverse and the transseries volume, 7~October 2026]\label{cmg:rem:transseries}
Compared statement by statement with the repository's transseries
volume~\cite{TSvol}; its symbols are not used here.
\begin{enumerate}\raggedright
\item \emph{The growth.} $\log U_n=\frac{n(n+1)}2\log n-\log n!+o(1)$ is not of the exponential--power form of \file{p0:def:model} (the leading term is $\frac12n^2\log n$), and the comparison function $G_D$ of~\eqref{cmg:eq:GD} keeps the exact Gamma function and exact sectors; no instance of the volume's model is claimed.
\item \emph{The threshold.} $N_*(y)$ is the volume's staircase (\file{p0:def:three-inverses}(1)) of the eventually strictly increasing sequence $U_n$. $G_D$ is not an interpolation of $\log U_n$ (it undershoots it by at most $B_D(n)$ at the integers), so Theorem~\ref{cmg:thm:inverse} is proved directly: its two-ceiling enclosure is an analogue of \file{p0:thm:staircase}(2), not an instance, with a one-sided error. The remark after it, that the single upper ceiling fails arbitrarily far out, is the phenomenon that \file{p0:thm:staircase}(2) and~(5) describe.
\item \emph{The initializer.} With $X=u^2$, the equation $H(u)=L$ reads $X\cdot\frac14\log X=L$, the volume's factorial core \file{p0:prop:factorial-core} with slope parameter $\frac14$ and shift $0$; its solution $X=\exp W_0(4L)=4L/W_0(4L)$ gives exactly~\eqref{cmg:eq:u}. So $u$ is the square root of an exact instance. Equivalently, in the logarithmic chart of $u$ the equation is $2\log u+\log\log u-\log2=\log L$, a monomial--logarithmic datum (\file{plt:def:lw-monomial-log-datum}) with exponent $2$, logarithmic power $1$, constant $-\log 2$ and no tail, so $u$ is also the $B=0$ case of \file{plt:thm:lw-template}. The correction in~\eqref{cmg:eq:initializer} comes from the Stirling terms of $F-H$, which are powers of $u^{-1}$, not of $1/\log u$; the write has not shown it to be an instance of either theorem, and the source proves it directly by one Newton step.
\end{enumerate}
\end{remark}

\section{Exact computations and the companion package}\label{cmg:sec:computations}
Finite tests support the proofs and catch implementation or transcription errors. They do not establish asymptotic validity. The public package separates integer/rational checks, symbolic cancellations, high-precision diagnostics, and deterministic artifact reproduction.

\subsection{Locally computed counts}
The following values fix the initial normalization:
\begin{center}
\begin{tabular}{r r}
\toprule
$n$&$U_n$\\\midrule
0&1\\
1&1\\
2&4\\
3&129\\
4&43968\\
5&254429900\\
6&30468670170912\\
7&91267244789189735259\\
8&8048575431238519331999571800\\
\bottomrule
\end{tabular}
\end{center}
Exact values through $n=24$ are supplied as decimal integer strings, together with comparison to the twelve displayed OEIS terms $n=0,\ldots,11$. These are locally computed data, not a copy of the official b-file. The b-file was not retrieved in the source checks underlying this report.

\begin{remark}[{[write]} The OEIS entry, 7~October 2026]\label{cmg:rem:oeis}
The write read A001425 in the internal format on 7~October 2026: revision \#43 (10 December 2025), ``Number of commutative groupoids with n elements.'', author N.~J.~A. Sloane, more terms by Christian G. Bower (1998), the cycle-type formula (dated 3 December 2003 in the entry's history), the asymptotic ``a(n) asymptotic to (n\textasciicircum binomial(n+1, 2))/n! = A023813(n)/A000142(n) \textasciitilde{} e\textasciicircum n*n\textasciicircum binomial(n, 2) / sqrt(2*pi*n)'' (unattributed), and a PARI program by Andrew Howroyd (10 December 2025), as Section~\ref{cmg:sec:sources} says. The entry links a b-file by Howroyd, ``Table of n, a(n) for n = 0..36'', which the source did not retrieve. \emph{Checks.} All 37 b-file terms equal the write's own Burnside sum, and the delivered counts for $n\le24$ (\file{data/counts.json}) equal the b-file; so the source's ``locally computed'' values agree with the official b-file. The entry contains no conjecture, so none is settled here; no OEIS edit is part of this work.
\end{remark}

For all $7{,}337$ positive-order partition types with $n\le24$, the exact checker constructs a representative permutation and compares its fixed-law count by a literal input-pair orbit walk with the trace/M\"obius method \eqref{cmg:eq:mobius} and the support product \eqref{cmg:eq:Fmu}. The same types test the orbit-loss inequalities, including fixed-point-free and zero-count examples. Summing with exact conjugacy-class multiplicities yields the Burnside integer; divisibility by $n!$ is checked explicitly.

\subsection{Coefficient and probability checks}
For all eleven moved types through defect $4$, the package computes $p_1,\ldots,p_5$ as exact rational numbers. The coefficient recurrence is independently checked by expanding the exact logarithm \eqref{cmg:eq:rawlog} and by exponentiating through partitions of the coefficient index. Symbolic checks also verify the pole and constant cancellation, the polynomial exponent, the Lambert-$W$ reduction, and the first inverse correction cancellation.

All $1+8+729=738$ labeled laws for $n=1,2,3$ are enumerated directly against every permutation. At $n=3$ the automorphism-order distribution is
\[
 \begin{array}{c|rrrr}
 |\Aut|&1&2&3&6\\\hline
 \text{labeled laws}&696&24&8&1
 \end{array}
\]
and the exact statistics include
\[
 \sum|\Aut|=774,\quad \sum X=27,\quad \sum Y=18,\quad
 \#Z=9,\quad \sum_ZX=3,\quad \sum_Z|\Aut|=30.
\]
Thus $U_3=774/6=129$, $P_3=11/243$, and $Q_3=13/129$. This verifies the change of measure in a range where larger groups genuinely occur; it does not assume that all nonrigid laws have only a transposition.

\subsection{Exact finite-tail checks, including odd $n$}
For $8\le n\le24$ and $0\le D\le\lfloor n/4\rfloor-1$, the package checks $R_D(n)$ exactly. Nonintegral exponents in $B_D(n)$ cannot be handled by silently rounding the bound upward and then claiming the original inequality was verified. Instead set
\begin{equation}\label{cmg:eq:rational-B}
 \underline B_D(n)=
 \frac{n^{-an+a(a+2)}}{1-n^{\lfloor3-n/2\rfloor}}
 +n^{\lfloor n-3n^2/16\rfloor}.
\end{equation}
For $n>1$, both replacements decrease the corresponding positive summands, hence $\underline B_D(n)\le B_D(n)$. The exact check $R_D(n)\le\underline B_D(n)$ is therefore stronger than the desired finite inequality at those test points. All comparisons use rational arithmetic; decimal ratios are diagnostics only. There are $62$ such $(D,n)$ pairs, including odd $n$.

\subsection{Reproduction and its limits}
The archive contains this manuscript, the compiled PDF, a standalone builder, portable Python code, locally generated data, and SHA-256 manifests. The README gives exact commands for a fresh extracted-ZIP build, an optimized-Python build, validation-only checks, and separate numerical diagnostics. Explicit exceptions, rather than Python \code{assert}, keep checks active under \code{python -O}.
\begin{writenote}[The shipped files, 7 October 2026]
In ProveIt the PDF, the two checksum manifests and \file{data/fixed_counts.json}
are not shipped; the report README lists the shipped files, gives the
command that regenerates \file{data/fixed_counts.json} from \file{code/check.py},
and explains that the builder and \file{code/test_guards.py} expect the
delivered layout and run only in a re-extracted archive.
\end{writenote}

The deterministic build fixes PDF metadata, stabilizes cross-references, and uses sorted ZIP members with fixed timestamps and attributes. With the matching recorded toolchain, the regenerated public members and whole ZIP are compared byte for byte. Hash manifests detect accidental alteration relative to the package's stated records; they do not authenticate an untrusted package or make arbitrary TeX/Python safe to execute. The malformed-input tests cover the documented interface and package checks, not every possible operating-system or hostile-input scenario.

The exact computations have a finite declared range; the proofs give the general statements. High-precision root and sector evaluations illustrate scales and test formula transcription. They are explicitly separate from the rigorous eventual inverse theorem and do not certify an effective inverse cutoff.

\section{Prior work and the boundary of the source review}\label{cmg:sec:sources}
The current A001425 entry \cite{oeis} supplies the exact cycle-type Burnside formula, the leading equivalent \eqref{cmg:eq:known-leading}, and a cycle-length implementation credited to Andrew Howroyd on 10 December 2025. These are established background for this report. The twelve displayed terms were read and compared exactly; the linked official b-file was not obtained. No data beyond the display are attributed to that file.

Freese's full 1990 paper \cite{freese} develops the relationship between labeled and unlabeled probabilities and proves asymptotic rigidity for unrestricted finite similarity types containing a higher-arity operation (with additional unary cases). It also discusses decay faster than inverse powers. This is substantial classical asymmetry background. Its unrestricted operation space is different from the commutativity-conditioned space used here, so its theorem is not invoked as an automatic proof for this model. The change-of-measure identity and the transposition/nontransposition event framework are reproduced explicitly.

Ture\v cek's full 2023 paper \cite{turecek} gives exact enumeration and cycle-index methods for unrestricted finite magmas, with higher-arity corrections to classical formulas. It provides context for the exact counting mechanism, rather than the commutative global remainder proved here. Burris's 2025 author exposition \cite{burris}, also read in full, presents unrestricted rigidity and asymptotics as a rewrite of parts of Freese. Neither unrestricted result is treated as a substitute for the commutative input-orbit analysis.

Harrison's two 1966 papers \cite{harrison,harrison-note} and Tamura's 1970 chapter \cite{tamura} are important historical references. Their full originals were not retrieved for this review; the Tamura scan linked from OEIS was inaccessible in the attempted retrievals. Their complete theorem scope is therefore not characterized here. Older Murskii originals appearing in the later bibliography were likewise not reviewed in full.
\begin{writenote}[7 October 2026]
The write did not read the Harrison or Tamura originals either; the comparison
stays open (Section~\ref{cmg:sec:questions}, question~6). The official b-file
that this section says was not obtained has been compared since: see
Remark~\ref{cmg:rem:oeis}.
\end{writenote}

The fixed-defect hierarchy developed here was not identified in the full texts actually read or in bounded exact-model searches of the available report repository and prior-artifact records. Such searches are not exhaustive. They do not establish worldwide originality, absence of an earlier proof, or complete nonoverlap with all historical material. The defensible contribution of this document is its self-contained derivation and reproducible verification of the stated refinements, with these source limits preserved.

\section{Further questions}\label{cmg:sec:questions}
\begin{enumerate}[leftmargin=*,itemsep=7pt]
\item \textbf{Subgroup-resolved rare-symmetry expansions.}
Determine further terms in $P_n$ and $Q_n$, organized by the permutation subgroups that stabilize a table. Individual automorphism sectors count elements, whereas nonrigidity is a union event. One needs intersections of equivariance conditions and their subgroup lattice, not a naive division by group orders. An exact subgroup-equivariant input/output orbit method may make the next terms accessible.
\item \textbf{Sharper global remainders.}
The word bound $n^{2d}$ and the estimate $q\le s(s+2)/4$ deliberately sacrifice precision for uniformity. Can the complete omitted sum be bounded by a constant or controlled polynomial times its first omitted sector, retaining a finite-$n$ statement over macroscopic supports? Such a result would improve practical inverse errors considerably.
\item \textbf{Growing defect and support.}
The amplitude disk $|t|<1/s$ and the finite-family arguments are used for fixed $D$. Establish uniform coefficient and remainder estimates when $D=D(n)$ grows. Identify how far the lexicographic hierarchy remains quantitatively useful before collective contributions of many types matter.
\item \textbf{Effective inverse cutoffs.}
Prove explicit thresholds beyond which $G_D'>0$ and $G_D(n+1)>G_D(n)+B_D(n)$, and bound the derivative infimum in \eqref{cmg:eq:eta} by a computable expression. A directed-rounding implementation could then turn the eventual enclosure into a certified numerical inverse. Integer-boundary ambiguity must still be retained.
\item \textbf{Related equational classes.}
Repeat the input-orbit analysis for idempotent commutative operations, a distinguished identity element, or other constraints whose fixed-point structure is tractable. Associativity is a much stronger global restriction and does not follow by changing the number of free table entries.
\item \textbf{Completing historical comparison.}
Retrieve and inspect the Harrison originals and Tamura chapter, and broaden exact-model searches. This may sharpen attribution and reveal earlier special cases without changing the internal validity of the proofs.
\end{enumerate}
\begin{writenote}[Added 7 October 2026, under Vladimir's standing rule of 4 October 2026]
Every open question of the source is in this list, with its sketch and what is
missing; the write adds from the non-claims: a quantitative version of the
eventual ordering (the ratio $T_{22}/T_3$ still exceeds one at $n=25$), and
uniform-in-$\mu$ amplitude bounds (only fixed-$\mu$ Cauchy bounds are proved).
No claim of the source was found false, and the OEIS entry contains no
conjecture.
\end{writenote}

\section*{Conclusion}
The leading count of commutative magmas comes from the identity permutation. Its symmetry excess has a finer structure: exact nonidentity cycle types supply convergent exponential sectors, and a global orbit-loss argument makes their fixed-defect truncations rigorous. The first transposition sector controls rare nonrigidity, with a factor of two between labeled and unlabeled sampling. Exact identity and sector terms also support inverse thresholds, provided the final integer rounding is kept as a two-ceiling enclosure. The companion computations make each finite claim reproducible while leaving the asymptotic arguments, the eventual qualifiers, and the source limitations visible.

\begin{thebibliography}{9}
\bibitem{oeis}
The OEIS Foundation Inc., \emph{The On-Line Encyclopedia of Integer Sequences}, entry A001425, ``Number of commutative groupoids with $n$ elements.''
\url{https://oeis.org/A001425}. Accessed 4 October 2026. Exact formula, displayed terms, and leading equivalent.

\bibitem{freese}
R. Freese, \emph{On the two kinds of probability in algebra}, Algebra Universalis \textbf{27} (1990), 70--79.
\href{https://doi.org/10.1007/BF01190254}{doi:10.1007/BF01190254}.
Author-hosted full text: \url{https://math.hawaii.edu/~ralph/Preprints/palgnew.pdf}.

\bibitem{turecek}
P. Ture\v cek, \emph{Counting Finite Magmas}, arXiv:2305.00269v2 (2023).
\url{https://arxiv.org/abs/2305.00269v2}.

\bibitem{burris}
S. Burris, \emph{Finite Magmas: Rigidity and Asymptotics}, author exposition, 13 May 2025, CountUpdate repository.
\href{https://github.com/snburris/CountUpdate/blob/main/count_magma_eqns/Rigidity\%20and\%20Asymptotics.pdf}{Author's PDF in CountUpdate}.

\bibitem{harrison}
M. A. Harrison, \emph{The number of isomorphism types of finite algebras}, Proceedings of the American Mathematical Society \textbf{17} (1966), 731--737.
\href{https://doi.org/10.1090/S0002-9939-1966-0200219-9}{doi:10.1090/S0002-9939-1966-0200219-9}. Original full text not retrieved for this review.

\bibitem{harrison-note}
M. A. Harrison, \emph{Note on the number of finite algebras}, Journal of Combinatorial Theory \textbf{1} (1966), 395--397.
\href{https://doi.org/10.1016/S0021-9800(66)80061-3}{doi:10.1016/S0021-9800(66)80061-3}. Original full text not retrieved for this review.

\bibitem{tamura}
T. Tamura, \emph{Some contributions of computation to semigroups and groupoids}, in J. Leech (ed.), \emph{Computational Problems in Abstract Algebra}, Pergamon, Oxford, 1970, pp. 229--261.
\href{https://oeis.org/A001329/a001329.pdf}{Annotated scan linked by OEIS}. Chapter not retrieved for this review.
\bibitem{TSvol}{\raggedright ProveIt, \emph{Transseries and inversion}, \file{Analysis/Transseries/docs/series-and-transseries/Transseries_And_Inversion/transseries_and_inversion.tex}; Part~0 (\file{p0:def:model}, \file{p0:prop:factorial-core}, \file{p0:def:three-inverses}, \file{p0:thm:staircase}) and the chapter on the Lambert template (\file{plt:def:lw-monomial-log-datum}, \file{plt:thm:lw-template}). [Added in the write.]\par}
\end{thebibliography}
\end{document}
